% GENERATED FILE. Edit canonical lessons, then run npm run generate:books.
% canonical-source-hash: fnv1a64:d21cf6aa4d3f2946
% canonical-lessons: TA-S137-letter-nya, TA-S138-letter-uu, TA-S139-letter-o, TA-S140-letter-ssa, TA-R112-letters-first-pass, TA-R112-letters-second-pass

\chapter{Letters From Days, Towns and Gladness}
\label{ch:ta-letters-from-days-towns-gladness}
\hlchaptermodality{112}

\begin{chapteropening}
\textbf{By the end of this chapter:} \emph{I can write the letters nya, uu, o and ssa, and find each one in a word I already say.}

Complete the last lesson of chapter 112: i can write the letters nya, uu, o and ssa, and find each one in a word I already say.
\end{chapteropening}

\section[ña]{\ta{ஞ} — the letter that opens \ta{ஞாயிறு}}

\label{lesson:TA-S137-letter-nya}



\begin{quote}
Before the new one: read \ta{◌ூ} and \ta{ஏ} aloud, and write each once.

Now say \textbf{\ta{ஞாயிறு}} (\emph{ñāyiṟu}), \textquotedblleft{}Sunday\textquotedblright{}. You have read it for a long time. This lesson takes one character out of it and puts it in your hand.
\end{quote}

\subsection*{Script you'll notice: \ta{ஞ}}
\textbf{\ta{ஞ}} — \emph{ña}.

It opens \textbf{\ta{ஞாயிறு}} (\emph{ñāyiṟu}), \textquotedblleft{}Sunday\textquotedblright{}, the last day in the week you learned. \ta{ஞ} is a rare letter: few everyday words use it, and Sunday is the one you will see most.

\subsection*{Writing: \ta{ஞ}}
\hlblockfigure{\detokenize{figures/TA-S137-letter-nya-filmstrip.pdf}}{How \ta{ஞ} is written, stroke by stroke}

\begin{itemize}
\raggedright
  \item \textbf{1.} start here: curve down into the compact left inner turn
  \item \textbf{2.} without lifting, circle the left inner loop and return upward, then lift
  \item \textbf{3.} put the pen down again and carry the long top bar straight to the right, then lift
  \item \textbf{4.} put the pen down again and curve left from the upper-right shoulder
  \item \textbf{5.} without lifting, descend the central upright, then lift
  \item \textbf{6.} put the pen down again and sweep around the broad outer-right curve
  \item \textbf{7.} without lifting, continue around the broad bottom bowl
  \item \textbf{8.} without lifting, return up the outer-left side — and only now lift
\end{itemize}

\textbf{Pen lifts: 3.}

\begin{quote}
Stroke order is one attested teaching order, not a national standard — Tamil handwriting is taught with school-to-school variation. Source: Sankaran Radhakrishnan, Tamil Script Learners Manual, Appendix I: Hand-movements, Frame 8, \ta{ஞ} (University of Texas at Austin), p. 194.
\end{quote}

\subsection*{Your turn}
\begin{itemize}
\raggedright
  \item \emph{Look:} at \textbf{\ta{ஞாயிறு}} and point at \ta{ஞ} inside it
  \item \emph{Trace it:} \ta{ஞ} three times, saying \emph{ña} as you finish each one
  \item \emph{Write it:} \ta{ஏ}, then \ta{ஞ}, from memory
  \item \emph{Read it:} \textbf{\ta{ஞாயிறு}} aloud once more
\end{itemize}

\begin{tcolorbox}[breakable,colback=teal!4,colframe=teal!35!black,title={Before you move on}]
Which character is this — \ta{ஞ}? (\textbf{ña}.) Which word did it come from? (\textbf{\ta{ஞாயிறு}}, \emph{ñāyiṟu}, \textquotedblleft{}Sunday\textquotedblright{}.)
\end{tcolorbox}
\section[ū]{\ta{ஊ} — the letter that opens \ta{ஊர்}}

\label{lesson:TA-S138-letter-uu}



\begin{quote}
Before the new one: read \ta{ஏ} and \ta{ஞ} aloud, and write each once.

Now say \textbf{\ta{ஊர்}} (\emph{ūr}), \textquotedblleft{}town, home place\textquotedblright{}. You have read it for a long time. This lesson takes one character out of it and puts it in your hand.
\end{quote}

\subsection*{Script you'll notice: \ta{ஊ}}
\textbf{\ta{ஊ}} — \emph{ū}.

It opens \textbf{\ta{ஊர்}} (\emph{ūr}), \textquotedblleft{}town\textquotedblright{}, the place a Tamil introduction asks you about. It is built on \textbf{\ta{உ}} (\emph{u}), which you already write, with a second part added to its right: short vowel, then more pen for the long one.

\subsection*{Writing: \ta{ஊ}}
\hlblockfigure{\detokenize{figures/TA-S138-letter-uu-filmstrip.pdf}}{How \ta{ஊ} is written, stroke by stroke}

\begin{itemize}
\raggedright
  \item \textbf{1.} start here: sweep outward around the compact upper spiral
  \item \textbf{2.} without lifting, descend through the broad outer curve and turn left onto the baseline
  \item \textbf{3.} without lifting, carry the long baseline straight to the right, then lift
  \item \textbf{4.} put the pen down again and curl around the added left loop
  \item \textbf{5.} without lifting, turn inward through the added loop
  \item \textbf{6.} without lifting, descend around the added inner curl, then lift
  \item \textbf{7.} put the pen down again and rise on the added adjoining upright
  \item \textbf{8.} without lifting, carry the added top bar right, then lift
  \item \textbf{9.} put the pen down again and descend the added separate right upright — and only now lift
\end{itemize}

\textbf{Pen lifts: 3.}

\begin{quote}
Stroke order is one attested teaching order, not a national standard — Tamil handwriting is taught with school-to-school variation. Source: Sankaran Radhakrishnan, Tamil Script Learners Manual, Module 17, \ta{ஊ} construction, with Appendix I: Hand-movements, Frames 17, 16, and 12 (University of Texas at Austin), pp. 195–196.
\end{quote}

\subsection*{Your turn}
\begin{itemize}
\raggedright
  \item \emph{Look:} at \textbf{\ta{ஊர்}} and point at \ta{ஊ} inside it
  \item \emph{Trace it:} \ta{ஊ} three times, saying \emph{ū} as you finish each one
  \item \emph{Write it:} \ta{ஞ}, then \ta{ஊ}, from memory
  \item \emph{Read it:} \textbf{\ta{ஊர்}} aloud once more
\end{itemize}

\begin{tcolorbox}[breakable,colback=teal!4,colframe=teal!35!black,title={Before you move on}]
Which character is this — \ta{ஊ}? (\textbf{ū}.) Which word did it come from? (\textbf{\ta{ஊர்}}, \emph{ūr}, \textquotedblleft{}town, home place\textquotedblright{}.)
\end{tcolorbox}
\section[o]{\ta{ஒ} — the letter that opens \ta{ஒரு}}

\label{lesson:TA-S139-letter-o}



\begin{quote}
Before the new one: read \ta{ஞ} and \ta{ஊ} aloud, and write each once.

Now say \textbf{\ta{ஒரு}} (\emph{oru}), \textquotedblleft{}a, one\textquotedblright{}. You have read it for a long time. This lesson takes one character out of it and puts it in your hand.
\end{quote}

\subsection*{Script you'll notice: \ta{ஒ}}
\textbf{\ta{ஒ}} — \emph{o}.

It opens \textbf{\ta{ஒரு}} (\emph{oru}), \textquotedblleft{}a, one\textquotedblright{}, the form \ta{ஒன்று} takes in front of a noun. You already read its long partner \textbf{\ta{ஓ}} (\emph{ō}). Put them side by side: the long one carries more at its foot.

\subsection*{Writing: \ta{ஒ}}
\hlblockfigure{\detokenize{figures/TA-S139-letter-o-filmstrip.pdf}}{How \ta{ஒ} is written, stroke by stroke}

\begin{itemize}
\raggedright
  \item \textbf{1.} start here: circle the small left loop and climb into the crown
  \item \textbf{2.} without lifting, sweep through the large right loop and curl inward, then lift
  \item \textbf{3.} put the pen down again and draw the separate lower bowl and return left — and only now lift
\end{itemize}

\textbf{Pen lifts: 1.}

\begin{quote}
Stroke order is one attested teaching order, not a national standard — Tamil handwriting is taught with school-to-school variation. Source: Sankaran Radhakrishnan, Tamil Script Learners Manual, Module 14, \ta{ஒ}, with Appendix I: Hand-movements, Frame 14 (University of Texas at Austin), p. 195.
\end{quote}

\subsection*{Your turn}
\begin{itemize}
\raggedright
  \item \emph{Look:} at \textbf{\ta{ஒரு}} and point at \ta{ஒ} inside it
  \item \emph{Trace it:} \ta{ஒ} three times, saying \emph{o} as you finish each one
  \item \emph{Write it:} \ta{ஊ}, then \ta{ஒ}, from memory
  \item \emph{Read it:} \textbf{\ta{ஒரு}} aloud once more
\end{itemize}

\begin{tcolorbox}[breakable,colback=teal!4,colframe=teal!35!black,title={Before you move on}]
Which character is this — \ta{ஒ}? (\textbf{o}.) Which word did it come from? (\textbf{\ta{ஒரு}}, \emph{oru}, \textquotedblleft{}a, one\textquotedblright{}.)
\end{tcolorbox}
\section[ṣa]{\ta{ஷ} — the letter inside \ta{சந்தோஷம்}}

\label{lesson:TA-S140-letter-ssa}



\begin{quote}
Before the new one: read \ta{ஊ} and \ta{ஒ} aloud, and write each once.

Now say \textbf{\ta{சந்தோஷம்}} (\emph{santōṣam}), \textquotedblleft{}gladness\textquotedblright{}. You have read it for a long time. This lesson takes one character out of it and puts it in your hand.
\end{quote}

\subsection*{Script you'll notice: \ta{ஷ}}
\textbf{\ta{ஷ}} — \emph{ṣa}.

It is inside \textbf{\ta{சந்தோஷம்}} (\emph{santōṣam}), \textquotedblleft{}gladness\textquotedblright{}. Like \ta{ஸ}, it is a \textbf{Grantha} letter, kept for words Tamil took from Sanskrit.

\subsection*{Writing: \ta{ஷ}}
\hlblockfigure{\detokenize{figures/TA-S140-letter-ssa-filmstrip.pdf}}{How \ta{ஷ} is written, stroke by stroke}

\begin{itemize}
\raggedright
  \item \textbf{1.} start here: curl around the inner-left loop and descend to the lower-left join, then lift
  \item \textbf{2.} put the pen down again and climb around the rounded outer-left body and descend to the lower join, then lift
  \item \textbf{3.} put the pen down again and draw the lower joining bar from the left body to the right, then lift
  \item \textbf{4.} put the pen down again and loop around the right body and finish with its descending tail — and only now lift
\end{itemize}

\textbf{Pen lifts: 3.}

\begin{quote}
Stroke order is one attested teaching order, not a national standard — Tamil handwriting is taught with school-to-school variation. Source: Ratnakar Narale, Learn Tamil Through English/Hindi, Reading and Writing Tamil Granthakshars, Third Tamil Granthakshar \ta{ஷ}, Sanskrit Hindi Research Institute, p. 13.
\end{quote}

\subsection*{Your turn}
\begin{itemize}
\raggedright
  \item \emph{Look:} at \textbf{\ta{சந்தோஷம்}} and point at \ta{ஷ} inside it
  \item \emph{Trace it:} \ta{ஷ} three times, saying \emph{ṣa} as you finish each one
  \item \emph{Write it:} \ta{ஒ}, then \ta{ஷ}, from memory
  \item \emph{Read it:} \textbf{\ta{சந்தோஷம்}} aloud once more
\end{itemize}

\begin{tcolorbox}[breakable,colback=teal!4,colframe=teal!35!black,title={Before you move on}]
Which character is this — \ta{ஷ}? (\textbf{ṣa}.) Which word did it come from? (\textbf{\ta{சந்தோஷம்}}, \emph{santōṣam}, \textquotedblleft{}gladness\textquotedblright{}.)
\end{tcolorbox}
\section[patinālu eḻuttu]{(fourteen letters) — Fourteen letters, fourteen words — a first pass}

\label{lesson:TA-R112-letters-first-pass}



\begin{quote}
Write the letter that opens \ta{இல்லை}, then the sign that ends it.
\end{quote}

\subsection*{Your turn}
\noindent
\begin{tabularx}{\linewidth}{@{}>{\raggedright\arraybackslash}X>{\raggedright\arraybackslash}X>{\raggedright\arraybackslash}X@{}}
\toprule
\textbf{letter} & \textbf{sound} & \textbf{the word it came from} \\
\midrule
\ta{இ} & \emph{i} & \textbf{\ta{இல்லை}} \emph{illai} \\
\ta{◌ை} & \emph{ai} & \textbf{\ta{இல்லை}} \emph{illai} \\
\ta{ஸ} & \emph{sa} & \textbf{\ta{நமஸ்காரம்}} \emph{namaskāram} \\
\ta{ங} & \emph{ṅa} & \textbf{\ta{நீங்கள்}} \emph{nīṅgaḷ} \\
\ta{ள} & \emph{ḷa} & \textbf{\ta{நீங்கள்}} \emph{nīṅgaḷ} \\
\ta{◌ீ} & \emph{ī} & \textbf{\ta{நீ}} \emph{nī} \\
\ta{◌ெ} & \emph{e} & \textbf{\ta{பெயர்}} \emph{peyar} \\
\ta{ஐ} & \emph{ai} & \textbf{\ta{ஐந்து}} \emph{aintu} \\
\ta{◌ூ} & \emph{ū} & \textbf{\ta{மூன்று}} \emph{mūṉṟu} \\
\ta{ஏ} & \emph{ē} & \textbf{\ta{ஏழு}} \emph{ēḻu} \\
\ta{ஞ} & \emph{ña} & \textbf{\ta{ஞாயிறு}} \emph{ñāyiṟu} \\
\ta{ஊ} & \emph{ū} & \textbf{\ta{ஊர்}} \emph{ūr} \\
\ta{ஒ} & \emph{o} & \textbf{\ta{ஒரு}} \emph{oru} \\
\ta{ஷ} & \emph{ṣa} & \textbf{\ta{சந்தோஷம்}} \emph{santōṣam} \\
\bottomrule
\end{tabularx}

\begin{itemize}
\raggedright
  \item \emph{Write it:} each letter in the table from memory, then check it
  \item \emph{Say it:} the word each one came from
  \item \emph{Write it:} the four vowel signs on their own: \ta{◌ை} \ta{◌ீ} \ta{◌ெ} \ta{◌ூ}
\end{itemize}

\begin{tcolorbox}[breakable,colback=teal!4,colframe=teal!35!black,title={Before you move on}]
Which two letters are Grantha, kept for borrowed words? (\textbf{\ta{ஸ}} and \textbf{\ta{ஷ}}.)
\end{tcolorbox}
\section[patinālu eḻuttu]{(fourteen letters) — From the word back to the letter — a second pass}

\label{lesson:TA-R112-letters-second-pass}



\begin{quote}
Say \emph{nīṅgaḷ}, then write the two new letters it holds.
\end{quote}

\subsection*{Your turn}
\noindent
\begin{tabularx}{\linewidth}{@{}>{\raggedright\arraybackslash}X>{\raggedright\arraybackslash}X>{\raggedright\arraybackslash}X@{}}
\toprule
\textbf{letter} & \textbf{sound} & \textbf{the word it came from} \\
\midrule
\ta{இ} & \emph{i} & \textbf{\ta{இல்லை}} \emph{illai} \\
\ta{◌ை} & \emph{ai} & \textbf{\ta{இல்லை}} \emph{illai} \\
\ta{ஸ} & \emph{sa} & \textbf{\ta{நமஸ்காரம்}} \emph{namaskāram} \\
\ta{ங} & \emph{ṅa} & \textbf{\ta{நீங்கள்}} \emph{nīṅgaḷ} \\
\ta{ள} & \emph{ḷa} & \textbf{\ta{நீங்கள்}} \emph{nīṅgaḷ} \\
\ta{◌ீ} & \emph{ī} & \textbf{\ta{நீ}} \emph{nī} \\
\ta{◌ெ} & \emph{e} & \textbf{\ta{பெயர்}} \emph{peyar} \\
\ta{ஐ} & \emph{ai} & \textbf{\ta{ஐந்து}} \emph{aintu} \\
\ta{◌ூ} & \emph{ū} & \textbf{\ta{மூன்று}} \emph{mūṉṟu} \\
\ta{ஏ} & \emph{ē} & \textbf{\ta{ஏழு}} \emph{ēḻu} \\
\ta{ஞ} & \emph{ña} & \textbf{\ta{ஞாயிறு}} \emph{ñāyiṟu} \\
\ta{ஊ} & \emph{ū} & \textbf{\ta{ஊர்}} \emph{ūr} \\
\ta{ஒ} & \emph{o} & \textbf{\ta{ஒரு}} \emph{oru} \\
\ta{ஷ} & \emph{ṣa} & \textbf{\ta{சந்தோஷம்}} \emph{santōṣam} \\
\bottomrule
\end{tabularx}

\begin{itemize}
\raggedright
  \item \emph{Say it:} each word in the table, then write the letter it gave you without looking
  \item \emph{Write it:} \ta{ஒ} beside \ta{ஓ}, and \ta{ஊ} beside \ta{உ}, and say how each pair differs
  \item \emph{Write it:} \ta{◌ெ} beside \ta{◌ே}, and \ta{◌ீ} beside \ta{◌ி}
\end{itemize}

\begin{tcolorbox}[breakable,colback=teal!4,colframe=teal!35!black,title={Before you move on}]
Which letters open a word, and which signs sit inside one? (\textbf{\ta{இ} \ta{ஐ} \ta{ஏ} \ta{ஊ} \ta{ஒ}} open words; \textbf{\ta{◌ை} \ta{◌ீ} \ta{◌ெ} \ta{◌ூ}} sit on a consonant.)
\end{tcolorbox}
